В первом примере массив состоит из $5$ элементов, изначально равных нулю. Рассмотрим состояние массива после каждой операции:

\begin{itemize}
\item после первой операции ($L_1 = 1$, $R_1 = 3$, $D_1 = 2$): $2\ 2\ 2\ 0\ 0$
\item после второй операции ($L_2 = 2$, $R_2 = 5$, $D_2 = -1$): $2\ 1\ 1\ -1\ -1$
\item после третьей операции ($L_3 = 4$, $R_3 = 4$, $D_3 = 5$): $2\ 1\ 1\ 4\ -1$
\end{itemize}

Итоговый массив после применения всех операций: <<2 1 1 4 -1>>.

Во втором примере массив состоит из единственного элемента, равного нулю. Единственная операция прибавляет $7$ к единственному элементу массива ($L_1 = R_1 = 1$, $D_1 = 7$), поэтому итоговый ответ: <<7>>.
